% page 103 du fac-simile (image p-102 du scan)
% bloc 165.1 x 237.7 mm ; marge gauche 36.9 mm
\pgc{15.240mm}{0.000mm}{
\l{\marge[77.774mm]{17.351mm}{\ornamento{12.687mm}{ornements/letroj/litero-D.png}}\cel{54}95}
\l{}
\l{}
\l{}
\l{}
\l{}
\l{d.- (\sou{Muziko}). Noto duesma di la \textquotedbl{}do\textquotedbl{}-gamo.}
\l{}
\l{da.- Indikas la individuo, objekto, kozo o kauzo qua igas,}
\l{\cel{3}facas od efektigas ulo.}
\l{}
\l{\sou{dafno}.- (\sou{Bot}.) Planto, genero de la familio \textquotedbl{}timelei\textquotedbl{}, analoga}
\l{\cel{3}a lauro.}
\l{}
\l{\sou{dago}.- Poniardo olima, di qua la lamo akuta povis penetrar tra}
\l{\cel{3}la lakuni di la kuraso o tra la mashkuti. - EFIS.}
\l{}
\l{\sou{\textquotedbl{}dahlia\textquotedbl{}}.-(\sou{Bot}.) Planto exotika, de la familio \textquotedbl{}kompozaji\textquotedbl{},}
\l{\cel{3}kultivata en nia gardeni pro lua flori di qui la kapituli}
\l{\cel{3}esas nereguloza, la kolori brilanta e diversa. - L.Dahlia.DEF.}
\l{}
\l{\sou{daimo}.- (\sou{Zool.})\cel{2}Genero di animalo, analoga a cervo, ma plu}
\l{\cel{3}mikra, e di qua la supra parto di la kornobranchi esas aplast-}
\l{\cel{3}ita e palmoida. - L.Cervus dama. - F.}
\l{}
\l{\sou{daktilo}.- Verso-pedo ek un silabo longa, sequata da du kurta. -}
\l{\cel{3}DEFIS.}
\l{}
\l{\sou{dalmatiko}.- (\sou{Liturgio}). Ornivo kirko-vestala, sorto di chazublo}
\l{\cel{3}sen-kruca, quan metas la diakono e la subdiakoni qui helpas la}
\l{\cel{3}sacerdoto. - DEFIRS.}
\l{}
\l{\sou{daltonismo}.- (\sou{Patol}.) Stando anormala di la vid-organo, qua}
\l{\cel{3}impedas dicernar ula kolori, precipue la kolori komplemen-}
\l{\cel{3}tala. - DEFIRS.}
\l{}
\l{\sou{damo}. - l. Muliero spozoza. II. Muliero vidva. III. En la karto-}
\l{\cel{3}ludo : figuro qua reprezentas rejino. IV. En la shako-ludo :}
\l{\cel{3}ta rejino qua unika darfas serchar e kaptar omna-direcione}
\l{\cel{3}ed irga-diste.\cel{2}V. En la piono-ludo : piono qua arivinte a}
\l{\cel{3}la lasta lineo di la piono-plako, duopligesas por distingar}
\l{\cel{3}lu de la ceteri, e darfas marchar e kaptar diagonale irga-}
\l{\cel{3}diste. - DEFIS.}
\l{}
\l{\sou{damasko}.- Stofo de silko, di qua la texuro prizentas flori,}
\l{\cel{3}desegnuri; stofo, linjo, qua prizentas ornamenti analoga.}
\l{\cel{3}- DEFIRS.}
\l{}
\l{\sou{damaskinar}. - (\sou{trans}.). Inkrustar oro, arjento (en peco de}
\l{\cel{3}fero o de stalo). - DEFIRS.}
\l{}
\l{\sou{damijano}.- Tre grosa botelo fera, gresa, terakota, di qua la}
\l{\cel{3}ventro esas larja, la kolo kurta, maxim-multa-kaze envelopita}
\l{\cel{3}ye ozier-plektajo,ed ansoza, uzata por transportar liquidi.-}
\l{\cel{3}EFIS.}
}
